\MathBookSourcePage{320}
\subsection{边值问题}
\label{sec:ch11-boundary-value-problems}
\setcounter{thmcount}{0}
\setcounter{equation}{0}

我们先引入一些记号. 约定在向量值算子、函数及其相应空间的符号下加单下划线, 在矩阵值函数和算子的符号下加双下划线.

设 $p$, $\chvec v=(v_1,v_2)^t$ 和 $\chten\tau=(\tau_{ij})_{1\leq i,j\leq2}$ 是二元函数. 定义
\[
\MathBookFormulaID{formula-ch11-gradient-curl-scalar}
\chvec{\operatorname{grad}}p=
\begin{pmatrix}\partial p/\partial x_1\\ \partial p/\partial x_2\end{pmatrix},
\qquad
\chvec{\operatorname{curl}}p=
\begin{pmatrix}\partial p/\partial x_2\\ -\partial p/\partial x_1\end{pmatrix},
\]
\[
\MathBookFormulaID{formula-ch11-divergence-rotation-vector}
\operatorname{div}\chvec v=\frac{\partial v_1}{\partial x_1}+\frac{\partial v_2}{\partial x_2},
\qquad
\operatorname{rot}\chvec v=-\frac{\partial v_1}{\partial x_2}+\frac{\partial v_2}{\partial x_1},
\]
\[
\MathBookFormulaID{formula-ch11-gradient-curl-vector}
\chten{\operatorname{grad}}\chvec v=
\begin{pmatrix}
\partial v_1/\partial x_1&\partial v_1/\partial x_2\\
\partial v_2/\partial x_1&\partial v_2/\partial x_2
\end{pmatrix},
\quad
\chten{\operatorname{curl}}\chvec v=
\begin{pmatrix}
\partial v_1/\partial x_2&-\partial v_1/\partial x_1\\
\partial v_2/\partial x_2&-\partial v_2/\partial x_1
\end{pmatrix},
\]
\[
\MathBookFormulaID{formula-ch11-symmetric-gradient}
\chten\epsilon(\chvec v)=\frac12\left(\chten{\operatorname{grad}}\chvec v+
(\chten{\operatorname{grad}}\chvec v)^t\right),
\]
以及
\MathBookSourcePage{321}
\[
\MathBookFormulaID{formula-ch11-divergence-matrix}
\chvec{\operatorname{div}}\chten\tau=
\begin{pmatrix}
\partial\tau_{11}/\partial x_1+\partial\tau_{12}/\partial x_2\\
\partial\tau_{21}/\partial x_1+\partial\tau_{22}/\partial x_2
\end{pmatrix}.
\]
还定义
\[
\MathBookFormulaID{formula-ch11-delta-chi-matrices}
\chten\delta=\begin{pmatrix}1&0\\0&1\end{pmatrix},
\qquad
\chten\chi=\begin{pmatrix}0&-1\\1&0\end{pmatrix},
\]
以及矩阵之间的如下内积
\[
\MathBookFormulaID{formula-ch11-matrix-inner-product}
\chten\sigma:\chten\tau=\sum_{i=1}^2\sum_{j=1}^2\sigma_{ij}\tau_{ij}.
\]
矩阵的迹定义为
\[
\MathBookFormulaID{formula-ch11-matrix-trace}
\operatorname{tr}(\chten\tau)=\chten\tau:\chten\delta=\tau_{11}+\tau_{22},
\]
并且有如下关系(参见习题~\ref{exr:ch11-identity-symmetric-gradient}):
\MBSourceItem{1}
\begin{equation}
\MathBookFormulaID{formula-ch11-symmetric-gradient-identity}
\label{eq:ch11-symmetric-gradient-identity}
\chten\epsilon(\chvec v)=\chten{\operatorname{grad}}\chvec v-
\frac12(\operatorname{rot}\chvec v)\chten\chi.
\end{equation}

考虑位形空间 $\Omega\subseteq\mathbb R^2$ 中的各向同性弹性材料. 令 $\chvec u(x)$ 为位移, $\chvec f(x)$ 为体力. 在线性弹性的静力理论中, $\chvec u$ 满足的方程为
\MBSourceItem{2}
\begin{equation}
\MathBookFormulaID{formula-ch11-static-elasticity-equation}
\label{eq:ch11-static-elasticity-equation}
-\chvec{\operatorname{div}}\chten\sigma(\chvec u)=\chvec f
\quad\text{in }\Omega,
\end{equation}
其中应力张量 $\chten\sigma(\chvec u)$ 定义为
\MBSourceItem{3}
\begin{equation}
\MathBookFormulaID{formula-ch11-stress-tensor}
\label{eq:ch11-stress-tensor}
\chten\sigma(\chvec u)=2\mu\chten\epsilon(\chvec u)
+\lambda\operatorname{tr}(\chten\epsilon(\chvec u))\chten\delta.
\end{equation}
正常数 $\mu$ 和 $\lambda$ 称为 Lamé 常数. 假设
$ (\mu,\lambda)\in[\mu_1,\mu_2]\times(0,\infty)$, 其中 $0<\mu_1<\mu_2$.

设 $\Gamma_1$ 和 $\Gamma_2$ 是 $\partial\Omega$ 的两个开子集, 满足
$\partial\Omega=\overline{\Gamma}_1\cup\overline{\Gamma}_2$ 且
$\Gamma_1\cap\Gamma_2=\varnothing$. 在 $\Gamma_1$ 上施加位移边界条件
\MBSourceItem{4}
\begin{equation}
\MathBookFormulaID{formula-ch11-displacement-boundary-condition}
\label{eq:ch11-displacement-boundary-condition}
\chvec u|_{\Gamma_1}=\chvec g,
\end{equation}
并在 $\Gamma_2$ 上施加牵引边界条件(其中 $\chvec\nu$ 是单位外法向量)
\MBSourceItem{5}
\begin{equation}
\MathBookFormulaID{formula-ch11-traction-boundary-condition}
\label{eq:ch11-traction-boundary-condition}
\left.\chten\sigma(\chvec u)\chvec\nu\right|_{\Gamma_2}=\chvec t.
\end{equation}
若 $\Gamma_1=\varnothing$(相应地, $\Gamma_2=\varnothing$), 则该边值问题称为纯牵引问题(相应地, 纯位移问题).

注意 $\chten\epsilon$ 有非平凡核. 令
\MBSourceItem{6}
\begin{equation}
\MathBookFormulaID{formula-ch11-rigid-motions-space}
\label{eq:ch11-rigid-motions-space}
\chvec{\mathrm{RM}}:=\left\{\chvec v:\chvec v=\chvec c+b(x_2,-x_1)^t,
\ \chvec c\in\mathbb R^2,\ b\in\mathbb R\right\}
\end{equation}
\MathBookSourcePage{322}
为无穷小刚性运动空间. 容易验证(参见习题~\ref{exr:ch11-rigid-motion-kernel})
\MBSourceItem{7}
\begin{equation}
\MathBookFormulaID{formula-ch11-rigid-motion-strain-zero}
\label{eq:ch11-rigid-motion-strain-zero}
\chten\epsilon(\chvec v)=0\qquad\forall\chvec v\in\chvec{\mathrm{RM}}.
\end{equation}
因此, $\chvec{\mathrm{RM}}$ 位于齐次纯牵引问题的核中, 而且要使纯牵引问题可解, $\chvec f$ 和 $\chvec t$ 必须满足某些约束.

Lebesgue 空间和 Sobolev 空间及其相应范数(以及在适当时的内积)按类似方式定义. 例如, 空间 $\chvec L^2(\Omega)$ 的内积为
\[
\MathBookFormulaID{formula-ch11-vector-l2-inner-product}
(\chvec\sigma,\chvec\tau)_{\chvec L^2(\Omega)}
:=\int_\Omega\chvec\sigma\cdot\chvec\tau\,dx.
\]
空间 $\chvec H^1(\Omega)$ 的内积为
\[
\MathBookFormulaID{formula-ch11-vector-h1-inner-product}
(\chvec u,\chvec v)_{\chvec H^1(\Omega)}
:=(\chten{\operatorname{grad}}\chvec u,\chten{\operatorname{grad}}\chvec v)_{\chten L^2(\Omega)}
+(\chvec u,\chvec v)_{\chvec L^2(\Omega)}.
\]
其他内积(包括最近使用的内积)与范数的定义留作习题.
